\scanpage{203}

Et on trouve donc que l'endomorphisme $F_*^{(\nu)}$ \deleted{de} $\omega$ n'est autre que l'\uline{hom trace} \inserted{$F_{X/S*}^{(\nu)}$ \deleted{\unsure{$=$}}}, associé à l'endom $F_{X/S}^{(\nu)}=(F_X)^\nu$ de $X/S$, [qui est transposé de l'endomorphisme $F_{X/S}^{(\nu)*}$ sur $\mathbb R^{n-i}f_*(\underline\Omega^{n-1}_{X/S})$, quand $X$ est propre et lisse sur $S$ --- ].

\deleted{Plaçons nous dans le cas} Notons que nous avons la suite exacte

\diagram{\begin{tikzcd}[ampersand replacement=\&,column sep=small]
0\arrow[r] \& \underline\omega\arrow[r]\arrow[d,equal] \& \mathcal H^1_{\mathrm{DR}}(X/S)\arrow[r]\&\deleted{\unsure{R}}\mathrm R^1 f_*(\underline O_{X/S})\\
\& \mathcal H^{1,0}_{\deleted{\unsure{X/S|\mathrm{DR}}}}(X/S) \& \&
\end{tikzcd}}

[avec un $0$ à droite dans les cas favorables), et que l'hom de Gysin \inserted{sur $\underline\omega$} associé à $F_{X/S}^{(\nu)}$ est compatible avec un hom de $=$ nom sur $\mathcal H^1_{\mathrm{DR}}(X/S)\longrightarrow\mathcal H^1_{\mathrm{DR}}(X^{p^\nu/S})$ \inserted{ss réserve $\unsure{\simeq|H}\ \mathcal H^1_{\mathrm{DR}}(X/S)^{(p^\nu)}$}, (qui \deleted{\unsure{sont|est}} des endom si $S=\mathrm{Spec}\,\mathbb F_{p^\nu}$). \deleted{Donc}

D'ailleurs, dans ce cas, \inserted{et si $X/S$ est propre,} \deleted{sauf erreur}, cet hom de Gysin induit zéro sur $\mathcal H^1_{\mathrm{DR}}(X/S)/\underline\omega\simeq\mathrm R^1 f_*(\underline O_{X/S})$, \deleted{il est peut-être \unsure{$=$|même|$\equiv$} nul sur les} (\unsure{regarder|raisonner|conjecture} par transposition), donc si $S=\mathrm{Spec}\,\unsure{\mathbb F_{p^\nu}|\#p^\nu}$, ses valeurs propres dans $\mathcal H^1_{\mathrm{DR}}(X/S)$ \deleted{et celles qui sont $\ne0$} sont aussi celles dans $\underline\omega$, et \deleted{\unsure{on}} \inserted{celles qui sont $\ne0$} sont aussi les valeurs propres \inserted{mod $p$} de Frobenius $\ell$-adique ($\ell\ne p$) \deleted{\unsure{sont cristallin|ou cristallin}}, ou cristallin qui sont des unités $p$-adiques. (Katz),

On voit donc que les $\omega$ tels que $\widetilde C\omega=\widetilde\omega$

% archive sheet number: 202
